% Smart Library & Fine Management System — CSE 3206 Sessional Report
% Compile with: xelatex report.tex   (or: lualatex report.tex)
% Run twice for correct TOC/figure numbering.
%
% Font note: this document requests Times New Roman via fontspec. If your
% machine does not have that font (common on Linux/Overleaf), the fallback
% TeX Gyre Termes is used — Termes is a Times-compatible font but is NOT
% Times New Roman. Do not claim otherwise in submissions.
\documentclass[12pt,a4paper]{article}

% ---------- fonts & encoding ----------
\usepackage{fontspec}
\IfFontExistsTF{Times New Roman}{
  \setmainfont{Times New Roman}
}{
  \setmainfont{TeX Gyre Termes} % substitute: Times-compatible, not Times New Roman
}

% ---------- layout & graphics ----------
\usepackage{amsmath}
\usepackage[top=1in,bottom=1in,left=1.25in,right=1.25in]{geometry}
\usepackage{graphicx}
\usepackage{float}
\usepackage{booktabs}
\usepackage{array}
\usepackage{longtable}
\usepackage{hyperref}
\hypersetup{colorlinks=true,linkcolor=black,citecolor=black,urlcolor=blue}
\usepackage{listings}
\usepackage{xcolor}
\usepackage{caption}
\usepackage{titlesec}
\usepackage{enumitem}
\usepackage{setspace}

\lstset{
  basicstyle=\ttfamily\footnotesize,
  breaklines=true,
  frame=single,
  backgroundcolor=\color{gray!8},
  keywordstyle=\color{blue!70!black},
  commentstyle=\color{green!40!black},
  showstringspaces=false,
  tabsize=2
}

\titleformat{\section}{\large\bfseries}{\thesection}{0.6em}{}
\titleformat{\subsection}{\normalsize\bfseries}{\thesubsection}{0.6em}{}
\titleformat{\subsubsection}{\normalsize\itshape}{\thesubsubsection}{0.6em}{}

\setlength{\parskip}{0.4em}
\setstretch{1.15}

\newcommand{\placeholder}[1]{\textbf{[PLACEHOLDER: #1]}}
\newcommand{\figpath}{figures/}

\title{\textbf{Smart Library \& Fine Management System}\\[0.5em]
\large Software Engineering Sessional Project Report}
\author{%
  Student Name: \placeholder{your name}\\
  Student ID: \placeholder{your ID}\\
  Supervisor/Instructor: \placeholder{instructor name}\\
  Department of Computer Science \& Engineering\\
  Rajshahi University of Engineering \& Technology (RUET)
}
\date{\placeholder{submission date}}

\begin{document}

\maketitle
\thispagestyle{empty}

\begin{abstract}
This report presents the design, implementation and verification of the
Smart Library \& Fine Management System, a full-stack web application
developed for a small academic library. The system follows the complete
Software Development Life Cycle using the Incremental Model and is built on
the MERN stack (MongoDB, Express, React, Node.js) with Tailwind CSS. Five
functional modules are implemented: member management, book catalog,
borrow/return with borrowing limits, overdue fine management, and audit with
in-application notifications. Three design patterns from three categories are
applied in production code: Singleton (shared database manager), Facade
(borrow/return workflow) and Strategy (interchangeable fine policies).
Verification combines Test-Driven Development with Jest, dependency mocking,
Supertest integration tests against an isolated in-memory MongoDB
(64 tests, 89.8\% statement coverage at the time of writing), a Selenium
WebDriver end-to-end test of the issue/return workflow, and a GitHub Actions
continuous-integration pipeline. Swagger/OpenAPI documentation, learned
independently, provides interactive API documentation at \texttt{/api-docs}.
The report documents honest limitations, including the absence of
authentication and of multi-document MongoDB transactions.

\vspace{1em}
\noindent\textbf{Keywords:} MERN stack, software engineering, design
patterns, Test-Driven Development, Selenium, continuous integration,
Swagger/OpenAPI, fine calculation.
\end{abstract}

\tableofcontents
\listoffigures
\newpage

% =========================================================
\section{Introduction}
% =========================================================
Libraries in academic institutions still manage borrowing records with
spreadsheets or manual registers, which are error-prone when calculating
overdue fines and tracking availability. This project builds a small but
complete web-based library information system that manages members, books,
borrowing and returning, overdue fines and an audit trail.

The project serves a dual purpose: (1) a practical product for a small
library, and (2) a vehicle to demonstrate Software Engineering techniques
required by CSE 3206: requirements engineering, process models, design
patterns, Test-Driven Development, mocking, integration testing, Selenium,
configuration management, continuous integration and an independently
learned tool (Swagger/OpenAPI).

\emph{Note on the assignment sheet:} the original specification contains an
inconsistency --- the heading states 10 marks while the marking scheme lists
20 marks. The student should confirm the official mark distribution with the
course instructor; this report does not assume either value.

% =========================================================
\section{Problem Statement}
% =========================================================
A small academic library must handle a growing set of members (students,
faculty and staff), a book catalog with limited copies, and an overdue fine
policy that differs by member category. Manual tracking leads to:
incorrect availability counts, forgotten due dates, inconsistent fine
calculations and no traceability of who did what. The problem is to design
and implement a maintainable, testable and extensible system that automates
these five concerns with persistent storage.

% =========================================================
\section{Objectives and Scope}
% =========================================================
\textbf{Objectives}
\begin{enumerate}[leftmargin=1.5em]
  \item Implement five functional modules with persistent MongoDB storage.
  \item Apply at least three design patterns from three categories in real
        workflows.
  \item Practice TDD, mocking and integration testing with measurable coverage.
  \item Use Git feature branches, pull requests and CI for configuration
        management.
  \item Document the API with Swagger/OpenAPI (independently learned tool).
  \item Produce a complete engineering report with UML and test evidence.
\end{enumerate}

\textbf{In scope:} role-labelled members (no authentication), issue/return
with limits, fine calculation, audit log, web UI, API documentation, CI.

\textbf{Out of scope:} authentication/authorization, online payments, email
or SMS notifications, barcode/RFID, multi-branch libraries, pagination.

% =========================================================
\section{Functional Requirements}
% =========================================================
The numbered functional requirements (FR-1.1 to FR-5.4) across the five
modules are listed in Table~\ref{tab:fr-summary} and fully specified in
\texttt{docs/requirements.md}. Each FR is covered by at least one automated
test.

\begin{table}[H]
\centering
\caption{Functional requirement summary}
\label{tab:fr-summary}
\begin{tabular}{>{\bfseries}p{2.2cm}p{10.5cm}}
\toprule
Module & Requirements \\
\midrule
Members & FR-1.1 register with unique ID; FR-1.2 reject duplicates (409);
FR-1.3 validate fields/role; FR-1.4 search and filter. \\
Books & FR-2.1 add with unique ISBN and copies; FR-2.2 reject duplicates and
invalid counts; FR-2.3 search by title/author/ISBN; FR-2.4
$0 \le$ available $\le$ total. \\
Borrow/return & FR-3.1 issue with $+14$-day due date; FR-3.2 limits
(Student 3, Faculty 5, Librarian 0); FR-3.3 validate member/book/stock;
FR-3.4 return restores stock (capped); FR-3.5 list active/returned. \\
Fines & FR-4.1 whole overdue days, never negative; FR-4.2 Student
10 BDT/day; FR-4.3 20 BDT/day after day 7; FR-4.4 Faculty 0 BDT;
FR-4.5 on-time return = 0; FR-4.6 overdue list with live fines. \\
Audit & FR-5.1 record issue/return/fine/overdue/member/book events;
FR-5.2 typed events with message and timestamp; FR-5.3 success/error
messages surfaced in the UI; FR-5.4 filterable history. \\
\bottomrule
\end{tabular}
\end{table}

\textbf{Fine policy boundary conditions (FR-4.2/4.3):}
let $d$ be whole overdue days. For students,
\[
\text{fine}(d) =
\begin{cases}
10d & 0 \le d \le 7\\
70 + 20(d-7) & d > 7
\end{cases}
\]
so $d=7 \Rightarrow 70$ BDT and $d=8 \Rightarrow 90$ BDT. Faculty:
$\text{fine}(d)=0\;\forall d$. A book returned on or before the due date has
$d=0$ and fine $0$.

% =========================================================
\section{Non-Functional Requirements}
% =========================================================
Eight non-functional requirements cover usability (one screen per module,
visible feedback), maintainability (small single-responsibility modules),
data integrity (unique IDs, atomic stock updates, non-negative fines),
reliability (centralised error handler, consistent JSON errors),
basic security (no credentials in code, input validation), performance
(sub-second responses for a small library), testability (DI + mocks +
isolated test DB) and extensibility (new fine policy = one new class).
Full text: \texttt{docs/requirements.md} \S 2.

% =========================================================
\section{User Stories and Use Cases}
% =========================================================
Nine user stories (e.g.\ ``As a librarian, I want to process a return that
automatically calculates any fine, so that I do not have to count overdue
days by hand'') and three fully written use cases (Issue a Book; Return a
Book and Calculate Fine; View Overdue Fines) with actors, preconditions,
main and alternative flows are documented in \texttt{docs/requirements.md}
\S 3--4.

% =========================================================
\section{SDLC Process Model}
% =========================================================
The \textbf{Incremental Model} was selected (Figure~\ref{fig:sdlc}). The
five modules are largely independent; each increment produces a working,
testable slice: (I1) requirements and design, (I2) members and books,
(I3) borrow/return, (I4) fines and audit, (I5) testing hardening, Swagger,
CI, Selenium and documentation. Incremental development matches a
single-developer sessional project: risks are isolated per increment, and a
defect in the fine policy cannot destabilise member registration because
they are separate increments and separate packages.

\begin{figure}[H]
\centering
\IfFileExists{\figpath sdlc-incremental.png}{
  \includegraphics[width=0.95\textwidth]{\figpath sdlc-incremental.png}
}{
  \fbox{\parbox{0.9\textwidth}{\centering\vspace{1em}
  \placeholder{render docs/diagrams/sdlc-incremental.mmd to
  figures/sdlc-incremental.png (see figures/README.md)}\vspace{1em}}}
}
\caption{Incremental development model used in this project
(source: \texttt{docs/diagrams/sdlc-incremental.mmd})}
\label{fig:sdlc}
\end{figure}

% =========================================================
\section{System Architecture}
% =========================================================
Three tiers: a React single-page application talks JSON/HTTP to an Express
REST API, which persists through Mongoose to MongoDB. Inside the server,
controllers are thin; the \textbf{BorrowFacade} coordinates multi-step
workflows; repositories isolate persistence; the \textbf{FineService}
selects a \textbf{Strategy} for each fine policy; \textbf{DatabaseManager}
(Singleton) owns the single connection manager.
Figure~\ref{fig:arch} shows the component view.

\begin{figure}[H]
\centering
\IfFileExists{\figpath architecture.png}{
  \includegraphics[width=0.9\textwidth]{\figpath architecture.png}
}{
  \fbox{\parbox{0.85\textwidth}{\centering\vspace{1em}
  \placeholder{create figures/architecture.png from the component description
  or the Mermaid source in docs/diagrams/}\vspace{1em}}}
}
\caption{Three-tier architecture with pattern locations}
\label{fig:arch}
\end{figure}

% =========================================================
\section{Module Descriptions}
% =========================================================
\textbf{M1 Members:} unique \texttt{memberId}, roles Student/Faculty/Librarian,
search/filter; implemented in \texttt{memberController.js} +
\texttt{MemberRepository} + React page \texttt{Members.jsx}.
\textbf{M2 Books:} unique ISBN, copy counters guarded by schema validation and
an atomic conditional update; \texttt{bookController.js}, \texttt{Books.jsx}.
\textbf{M3 Borrow/return:} \texttt{BorrowFacade.issueBook/returnBook};
14-day loans; per-role limits; stock updates; \texttt{BorrowReturn.jsx}.
\textbf{M4 Fines:} strategy-based calculation, live overdue endpoint;
\texttt{Fines.jsx}.
\textbf{M5 Audit/notifications:} \texttt{AuditService} persists typed events;
\texttt{ToastContext.jsx} shows in-app success/error messages;
\texttt{Audit.jsx} renders the history.

% =========================================================
\section{UI Prototype}
% =========================================================
A low-fidelity wireframe was produced before implementation
(\texttt{docs/diagrams/ui-wireframe.md}, Figure~\ref{fig:proto}). It planned
the sidebar layout, dashboard cards, member/book forms and tables, the
issue/return workflow, the fines policy cards and the audit list. The
prototype is a planning artefact and is not presented as evidence of
functioning software.

\begin{figure}[H]
\centering
\IfFileExists{\figpath ui-wireframe.png}{
  \includegraphics[width=0.9\textwidth]{\figpath ui-wireframe.png}
}{
  \fbox{\parbox{0.85\textwidth}{\centering\vspace{1em}
  \placeholder{optional: render the wireframe sketch to
  figures/ui-wireframe.png}\vspace{1em}}}
}
\caption{UI wireframe (planning prototype)}
\label{fig:proto}
\end{figure}

% =========================================================
\section{MongoDB Data Model}
% =========================================================
Four collections: \texttt{Member}, \texttt{Book}, \texttt{BorrowRecord} and
\texttt{AuditEvent} (Table~\ref{tab:schema}). MongoDB is document-oriented:
there are no foreign keys. \texttt{BorrowRecord} stores the
\texttt{\_id} of the member and book as \emph{logical references}, resolved
in the API with Mongoose \texttt{populate()} when names are needed for the
UI. Indexes cover unique \texttt{memberId}/ISBN, text search on
title/author, and the overdue query (\texttt{status, dueDate}).

\begin{table}[H]
\centering
\caption{MongoDB collections (Mongoose schemas)}
\label{tab:schema}
\small
\begin{tabular}{>{\bfseries}p{2.6cm}p{9.8cm}}
\toprule
Collection & Fields \\
\midrule
Member & \_id, memberId (unique), name, email, role (Student/Faculty/Librarian), createdAt \\
Book & \_id, isbn (unique), title, author, totalCopies ($\ge 1$),
availableCopies ($0\ldots$totalCopies), timestamps \\
BorrowRecord & \_id, memberId $\rightarrow$ Member, bookId $\rightarrow$ Book,
borrowedAt, dueDate, returnedAt, status (active/returned), fineAmount ($\ge 0$) \\
AuditEvent & \_id, eventType (enum of 6 types), memberId?, borrowRecordId?,
message, createdAt (index, descending) \\
\bottomrule
\end{tabular}
\end{table}

\begin{figure}[H]
\centering
\IfFileExists{\figpath mongodb-relationships.png}{
  \includegraphics[width=0.95\textwidth]{\figpath mongodb-relationships.png}
}{
  \fbox{\parbox{0.9\textwidth}{\centering\vspace{1em}
  \placeholder{render docs/diagrams/mongodb-relationships.mmd to
  figures/mongodb-relationships.png}\vspace{1em}}}
}
\caption{Document relationships (logical references, not relational FKs)}
\label{fig:er}
\end{figure}

% =========================================================
\section{REST API}
% =========================================================
Thirteen endpoints (Table~\ref{tab:api}) are implemented, validated and
documented in the OpenAPI file \texttt{server/docs/openapi.yaml}. Errors use
a consistent body \verb|{"success": false, "message": "..."}| via one
centralised error handler; status codes: 400 validation, 403 role, 404 not
found, 409 conflict, 500 unexpected.

\begin{table}[H]
\centering
\caption{REST API endpoints}
\label{tab:api}
\small
\begin{tabular}{llp{6.8cm}}
\toprule
Method & URL & Purpose \\
\midrule
POST & /api/members & Register member (409 on duplicate ID) \\
GET & /api/members?q=\&role= & List/search members \\
GET & /api/members/:id & Member details \\
POST & /api/books & Add book (409 on duplicate ISBN) \\
GET & /api/books?q= & Search by title/author/ISBN \\
GET & /api/books/:id & Book details \\
POST & /api/borrowings/issue & Issue \{memberId, bookId\} (Facade) \\
POST & /api/borrowings/return/:id & Return + fine calculation \\
GET & /api/borrowings?status= & Active/returned records \\
GET & /api/fines/overdue & Overdue list with live fines \\
GET & /api/fines/preview/:id & Fine preview (read-only) \\
GET & /api/audits?eventType= & Audit history \\
GET & /api/stats & Dashboard statistics \\
\bottomrule
\end{tabular}
\end{table}

% =========================================================
\section{Design Patterns}
% =========================================================

\subsection{Singleton — DatabaseManager}
\textbf{Category: creational.} \textbf{Problem:} the application must not
create a new MongoDB connection manager per request; initialisation and
errors must be handled in one place. \textbf{Location:}
\texttt{server/src/database/DatabaseManager.js}. \texttt{getInstance()}
returns the single instance held in a static field; \texttt{connect()}
reuses an existing connection. \textbf{Benefit:} one owner of the connection
lifecycle, predictable shutdown in tests. \textbf{Honest note:} the MongoDB
driver already pools connections — Singleton does not mean one socket and
does not automatically improve throughput. \textbf{Rejected alternative:}
Prototype — cloning a database manager is meaningless here; there is no
state worth duplicating and duplicated managers would fragment the
connection lifecycle.

\begin{lstlisting}[caption={Singleton core (excerpt)},label=lst:singleton]
class DatabaseManager {
  static instance = null;
  static getInstance() {
    if (DatabaseManager.instance === null) {
      DatabaseManager.instance = new DatabaseManager();
    }
    return DatabaseManager.instance;
  }
  async connect(uri) { /* reuse or establish connection */ }
}
\end{lstlisting}

\subsection{Facade — BorrowFacade}
\textbf{Category: structural.} \textbf{Problem:} issuing a book touches
member lookup, role/limit checks, atomic stock decrement, record creation,
due-date assignment and audit logging; returning adds fine calculation.
Without a facade, controllers would duplicate and mix this workflow.
\textbf{Location:} \texttt{server/src/facades/BorrowFacade.js}. The
controller calls exactly \texttt{borrowFacade.issueBook(memberId, bookId)}
and \texttt{returnBook(id)}. Persistence and fine calculation are delegated
to injected repositories and services (kept small on purpose).
\textbf{Benefit:} thin controllers, one readable workflow, unit-testable
with mocks. \textbf{Rejected alternative:} Adapter — Adapter solves
\emph{interface incompatibility} with a third-party API; here all
collaborators are our own classes and the need is \emph{orchestration}, not
interface conversion.

\subsection{Strategy — Fine Calculation}
\textbf{Category: behavioural.} \textbf{Problem:} fine rules differ by
member category and are expected to change; burying \texttt{if role == ...}
in a controller would scatter the policy.
\textbf{Location:} \texttt{server/src/strategies/} —
\texttt{StudentFineStrategy} (10 BDT/day), \texttt{FacultyFineStrategy}
(0), \texttt{OverdueHeavyFineStrategy} (decorator: first 7 days at the base
rate, then 20 BDT/day). \texttt{FineService.getStrategyForRole()} selects
and composes them behind the shared \texttt{calculateFine(days)} interface.
\textbf{Benefit:} substitutable policies; a new policy (e.g.\ Staff) is a
new class plus one switch case — existing strategies remain untouched
(open/closed). \textbf{Rejected alternative:} State — State models one
object whose behaviour changes with its \emph{internal state}; here we have
a family of interchangeable algorithms chosen by an \emph{external}
attribute (role), which is Strategy's definition.

\begin{lstlisting}[caption={Strategy selection (excerpt)},label=lst:strategy]
getStrategyForRole(role) {
  switch (role) {
    case 'Student':
      return new OverdueHeavyFineStrategy(
        new StudentFineStrategy(10), 7, 20);
    case 'Faculty':
    case 'Librarian':
      return new FacultyFineStrategy();
    default:
      throw new Error(`Unknown member role: ${role}`);
  }
}
\end{lstlisting}

\textbf{Test evidence:} \texttt{tests/unit/fineStrategy.test.js} verifies
0/1/5/7/8/10/14-day boundaries and invalid inputs;
\texttt{tests/unit/borrowFacade.test.js} verifies facade coordination with
fully mocked dependencies.

% =========================================================
\section{UML Diagrams}
% =========================================================
Figure~\ref{fig:class} shows the class diagram containing the actual classes
of the codebase; Figures~\ref{fig:seq-issue} and \ref{fig:seq-return} show
the issue and return+fine sequences. Mermaid sources are in
\texttt{docs/diagrams/}.

\begin{figure}[H]
\centering
\IfFileExists{\figpath class-diagram.png}{
  \includegraphics[width=0.98\textwidth]{\figpath class-diagram.png}
}{
  \fbox{\parbox{0.9\textwidth}{\centering\vspace{1em}
  \placeholder{render docs/diagrams/class-diagram.mmd to
  figures/class-diagram.png}\vspace{1em}}}
}
\caption{UML class diagram (actual implemented classes)}
\label{fig:class}
\end{figure}

\begin{figure}[H]
\centering
\IfFileExists{\figpath sequence-issue.png}{
  \includegraphics[width=0.98\textwidth]{\figpath sequence-issue.png}
}{
  \fbox{\parbox{0.9\textwidth}{\centering\vspace{1em}
  \placeholder{render sequence-issue.mmd}\vspace{1em}}}
}
\caption{Sequence diagram: issuing a book via BorrowFacade}
\label{fig:seq-issue}
\end{figure}

\begin{figure}[H]
\centering
\IfFileExists{\figpath sequence-return.png}{
  \includegraphics[width=0.98\textwidth]{\figpath sequence-return.png}
}{
  \fbox{\parbox{0.9\textwidth}{\centering\vspace{1em}
  \placeholder{render sequence-return.mmd}\vspace{1em}}}
}
\caption{Sequence diagram: return with Strategy-based fine calculation}
\label{fig:seq-return}
\end{figure}

A use-case diagram is shown in Figure~\ref{fig:usecase}; sources:
\texttt{docs/diagrams/use-case.mmd} (UML notation) and
\texttt{use-case-graph.mmd} (renderable variant).

\begin{figure}[H]
\centering
\IfFileExists{\figpath use-case.png}{
  \includegraphics[width=0.95\textwidth]{\figpath use-case.png}
}{
  \fbox{\parbox{0.9\textwidth}{\centering\vspace{1em}
  \placeholder{render use-case.mmd / use-case-graph.mmd to
  figures/use-case.png}\vspace{1em}}}
}
\caption{Use-case diagram (actors and included sub-functions)}
\label{fig:usecase}
\end{figure}

% =========================================================
\section{Implementation Overview}
% =========================================================
\textbf{Backend.} Node.js 22 + Express 4. \texttt{app.js} performs manual
dependency injection: repositories $\rightarrow$ services $\rightarrow$
facade $\rightarrow$ controllers $\rightarrow$ routes, so tests can replace
any collaborator. Request validation uses a small custom middleware
(\texttt{validate.js}); errors flow to \texttt{errorHandler.js}.
Stock safety: \texttt{decrementAvailable} uses a conditional atomic update
\texttt{\{availableCopies: \{\$gt: 0\}\} with \$inc: -1}, so two concurrent
issues cannot both take the last copy. Multi-document transactions are
\emph{not} used (a local standalone MongoDB does not support them without a
replica set); the ordered facade workflow plus the atomic decrement is the
documented, honest alternative — it is safe at this scale but is not a
transactional guarantee.

\textbf{Frontend.} React 18 + Vite 5 + Tailwind CSS 3.4, React Router pages
(Dashboard, Members, Books, Issue/Return, Fines, Audit), Axios client, and a
\texttt{ToastContext} for in-app success/error notifications. Pages show
loading, empty and error states and validate required fields before
submitting.

% =========================================================
\section{TDD Workflow}
% =========================================================
Fine calculation was developed test-first. The recorded cycle
(\texttt{docs/tdd-workflow.md}, logs in \texttt{docs/evidence/}): with the
implementation in its unimplemented state, the Student strategy tests failed
(\texttt{expected 50, received 0}; 2 failed / 3 passed) — \textbf{Red};
after implementing \texttt{return days * perDayRate}, all 5 focused tests
passed — \textbf{Green}; extraction of the heavy-overdue decorator and
service wiring with the suite still green — \textbf{Refactor}. TDD was
chosen because fine rules are precise business logic whose silent breakage
is expensive; boundary tests (day 7 vs.\ day 8) force the edge cases to
exist.

% =========================================================
\section{Mocking and Integration Testing}
% =========================================================
\textbf{Mocking.} \texttt{tests/unit/borrowFacade.test.js} replaces
\texttt{memberRepository}, \texttt{bookRepository},
\texttt{borrowRecordRepository}, \texttt{fineService} and
\texttt{auditService} with \texttt{jest.fn()} objects. Assertions verify,
for example, that exactly one \texttt{BOOK\_ISSUED} audit event is recorded
on success and \emph{none} on failure paths — without any database.

\textbf{Integration.} \texttt{tests/integration/api.test.js} (28 tests)
drives the real Express app with Supertest against an isolated
MongoDB produced by \texttt{mongodb-memory-server} (or
\texttt{MONGODB\_URI\_TEST} if provided), with full cleanup between tests —
development and Atlas data cannot be touched. Workflows covered: member and
book creation, duplicates, validation failures, issue, unavailable book,
borrow limit, on-time return (fine 0), overdue returns (50 BDT at 5 days;
90 BDT at day 8 for students; 0 for faculty), fine endpoints, audits and
dashboard stats.

\textbf{Coverage (actual run, 64 tests, all passing):}
statements 89.77\%, branches 73.33\%, functions 91.01\%, lines 90.88\%
(\texttt{npm run test:coverage}; full table in Appendix~\ref{app:coverage}).
Branch coverage is lowest in the centralised error handler's rarely-hit
branches — identified as an improvement target. Coverage measures execution,
not correctness; it is reported together with the assertions' intent.

% =========================================================
\section{Selenium Testing}
% =========================================================
\texttt{e2e/tests/issueAndReturn.e2e.js} drives Chrome through the built
React application: open dashboard $\rightarrow$ add book $\rightarrow$
register member $\rightarrow$ issue the book via the UI selects
$\rightarrow$ assert the success toast (``Due date'') and the active
borrowing row $\rightarrow$ click Return $\rightarrow$ assert ``No fine''.
Stable \texttt{data-testid} attributes are used as selectors. The
orchestrator \texttt{e2e/run-e2e.js} starts the backend in in-memory demo
mode, serves the frontend with \texttt{vite preview} on port 4173, waits for
health, runs the test and shuts down. \textbf{Recorded result on the
development machine: E2E TEST PASSED (6/6 steps)}; output in
Appendix~\ref{app:e2e}. Selenium runs locally and in the CI job for
ubuntu-latest (Chrome preinstalled, Selenium Manager fetches ChromeDriver).

% =========================================================
\section{Git/GitHub Configuration Management}
% =========================================================
Branch naming convention and full command recipes are in
\texttt{docs/git-workflow.md}: \texttt{feature/member-management},
\texttt{feature/book-catalog}, \texttt{feature/borrow-return},
\texttt{feature/fine-strategy}, \texttt{test/unit-and-integration},
\texttt{test/selenium-e2e}, \texttt{docs/swagger-openapi},
\texttt{ci/github-actions}. Commit messages follow a short imperative
summary with a scope prefix (e.g.\ \texttt{feat: add heavy-overdue fine
strategy with boundary tests}).

\placeholder{Append Git evidence produced on YOUR repository:}
\texttt{git log --oneline --graph --all} screenshot (Appendix~\ref{app:git}),
contribution of your own commits, and the README as delivered in this
repository.

% =========================================================
\section{Pull-Request Review Evidence}
% =========================================================
The report requires at least one accepted PR and one rejected/revised PR
with reviews. \texttt{docs/git-workflow.md} \S 3 contains the review
checklist used (correctness, tests, readability, error handling, secrets,
docs, scope).

\placeholder{Paste real PR evidence: accepted PR link/screenshot with review
comments; rejected or revised PR link/screenshot with the change requests
and your follow-up commits (Appendix~\ref{app:pr}). Do not fabricate reviews.}

% =========================================================
\section{Merge-Conflict Resolution Evidence}
% =========================================================
A safe, reproducible conflict exercise (two branches editing the same README
line, merge, resolve markers, verify) is documented in
\texttt{docs/git-workflow.md} \S 4.

\placeholder{Include screenshots of the conflicted file with conflict
markers and of the resolved commit (Appendix~\ref{app:conflict}).}

% =========================================================
\section{Continuous Integration (GitHub Actions)}
% =========================================================
\texttt{.github/workflows/ci.yml} runs on pushes and pull requests to
\texttt{main} with three jobs: \textbf{backend-tests} (checkout, Node 20,
\texttt{npm ci} from the lockfile, \texttt{npm run test:coverage} with the
in-memory MongoDB — no private credentials required, coverage uploaded as an
artifact), \textbf{frontend-build} (\texttt{npm ci} + \texttt{vite build})
and \textbf{selenium-e2e} (build + preview + browser test on the runner).
Failures are diagnosed by opening the Actions run $\rightarrow$ failed job
$\rightarrow$ step log.

\placeholder{A green checkmark on GitHub only exists after you push and the
runners succeed — include a screenshot of the successful run in
Appendix~\ref{app:ci} and state the run number honestly.}

% =========================================================
\section{Swagger/OpenAPI — Self-Directed Tool}
% =========================================================
\textbf{What was learned and why.} API documentation that stays usable
after the developer leaves the room. \textbf{Application.}
\texttt{server/docs/openapi.yaml} (OpenAPI 3.0) documents every implemented
endpoint with request/response schemas and examples; \texttt{swagger-ui-express}
serves interactive docs at \texttt{/api-docs}; the raw spec is exposed at
\texttt{/api/openapi.json}. A live check on the development machine returned
HTTP 200 for both (Appendix~\ref{app:swagger}).

\textbf{Demonstration.} Open \texttt{http://localhost:5000/api-docs},
expand \texttt{POST /api/books}, press \emph{Try it out}, execute a sample
book creation, and observe the 201 response — the same call the React form
makes.

\textbf{Benefits vs.\ alternatives.} Compared with a README: structured,
versionable, and rendered as interactive docs instead of prose that drifts.
Compared with a Postman collection: shareable as a URL/spec, diff-able in
Git, and vendor-neutral. \textbf{Limitations:} the spec is maintained by
hand (no automatic generation from Express code), response schemas are
summarised rather than exhaustive, and authentication schemes are not
documented because the system has none.

\textbf{Learning process (personalise before submission):} resources used
were the OpenAPI specification (\texttt{swagger.io/specification}) and the
swagger-ui-express README; \placeholder{describe one difficulty you hit
(e.g.\ YAML indentation) and how you resolved it}.

% =========================================================
\section{Learning Log}
% =========================================================
The learning-log template with columns (date, task, resource, learning,
challenge, solution, evidence, time) plus recommended official resources for
React, Tailwind, Node/Express, MongoDB/Mongoose, Jest, Supertest, Selenium,
Git/GitHub, Actions, OpenAPI and LaTeX, and a resource-evaluation method, is
provided in \texttt{docs/learning-log.md}. Sample rows are explicitly marked
\textbf{[EXAMPLE]} and must be replaced with your real sessions.

\placeholder{Attach your completed learning log (Appendix~\ref{app:log}).}

% =========================================================
\section{Reflection Essay (500--700 words)}
% =========================================================
% Reflection essay — 500-700 words.
% NOTE: items marked [PERSONAL: ...] MUST be rewritten from your real
% experience (or deleted) before submission. Do not submit with markers.
When I started this project, the goal was not to build the biggest library
system in the world, but to walk through a complete software engineering
process on a system small enough to understand end-to-end. The Smart Library
\& Fine Management System uses the MERN stack --- React with Tailwind CSS on
the front, Express and Mongoose with MongoDB behind a REST API --- because it
let me keep one language (JavaScript) across the whole project and focus on
engineering practice rather than fighting tooling.

For the process model I chose the Incremental Model. The five modules ---
members, books, borrow/return, fines and audit --- are naturally separable,
so I could finish and test each slice before starting the next. This suited
a single-developer sessional project: if something went wrong in the fine
policy, member registration was already stable and demonstrable.
[PERSONAL: describe one real problem you hit in a specific increment and how
incremental delivery helped.]

Three design patterns did real work in the code. The \textbf{Singleton}
\texttt{DatabaseManager} gives the process one shared connection manager, so
the connection lifecycle has a single owner; I rejected Prototype because
cloning a database manager would be meaningless --- there is nothing useful
to copy. The \textbf{Facade} \texttt{BorrowFacade} hides the multi-step
issue/return workflow (member lookup, borrow limit, stock update, due date,
fine calculation, audit logging) behind two readable methods; I rejected
Adapter because no incompatible third-party interface needed translating.
The \textbf{Strategy} pattern implements the fine policies --- 10~BDT/day
for students, zero for faculty, and a heavy-overdue decorator charging
20~BDT/day after the first seven days. I rejected State because the fine
rules are interchangeable algorithms selected by role, not an object whose
behaviour changes with internal state.

Test-Driven Development was applied most strictly to fine calculation. I
recorded a real Red--Green cycle: the tests failed against the unimplemented
strategy (``five overdue days $\rightarrow$ 50~BDT'' expected 50, received
0), passed once the ten-taka-per-day rule was implemented, and the boundary
cases (day~7~$=$~70, day~8~$=$~90) were locked in with further tests.
Mocking mattered in the facade unit tests: repositories, the fine service
and the audit service are all Jest mocks, so the tests verify coordination
logic without a database and run in milliseconds. Integration tests with
Supertest cover the real HTTP path against an isolated in-memory MongoDB,
which catches wiring mistakes mocks cannot see. The Selenium test drives the
actual React UI in Chrome --- add book, add member, issue, return --- and
passed on my machine.
[PERSONAL: describe one real bug that Selenium or the integration tests
found that unit tests missed.]

Continuous integration through GitHub Actions re-runs tests, coverage and
the frontend build on every push, which turns ``it works on my machine''
into a shared, visible fact. As the self-directed tool I learned
\textbf{Swagger/OpenAPI}: one YAML file now generates interactive
documentation at \texttt{/api-docs}, and I can test any endpoint from the
browser. I chose it over a README because a specification is structured,
versionable and much harder to let drift out of date.
[PERSONAL: mention which Swagger resources you actually used and one real
difficulty, e.g.\ YAML indentation.]

The hardest part was [PERSONAL: replace with your real challenge --- e.g.\
understanding atomic MongoDB updates, React select state in Selenium, or
wiring dependency injection]. I resolved it by [PERSONAL: your real solution
steps]. The learning resources I relied on were mostly official
documentation (Mongoose, Jest, Express, Tailwind), which I evaluated by
whether they matched the project's library versions and whether I could
apply them within one sitting.

The system has clear limitations: no authentication, no payment flow, regex
search instead of Atlas Search, and no multi-document transactions --- the
issue workflow instead uses an atomic conditional stock update, which I
documented honestly rather than overclaiming. Next I want to learn proper
authentication (JWT with role checks) and containerised deployment, because
the realistic next step for this application is running it safely for real
users rather than adding more features.

[FINAL STEP: after your viva rehearsals and real test runs, re-read the
essay aloud, keep it between 500 and 700 words, and remove every bracketed
marker.]


% =========================================================
\section{Limitations}
% =========================================================
\begin{enumerate}[leftmargin=1.5em]
  \item No authentication/authorization; roles are labels only.
  \item Fines are computed but not paid or waived.
  \item In-memory demo mode (\texttt{start:memory}) is non-persistent by design.
  \item Regex-based search instead of Atlas Search; no pagination.
  \item No multi-document transactions; atomic conditional stock update is
        the documented alternative (safe at small scale, not a full
        transactional guarantee).
  \item The OpenAPI spec is hand-maintained and can drift if endpoints change
        without updating the YAML.
  \item CI status on GitHub must be evidenced by an actual successful run on
        your repository.
\end{enumerate}

% =========================================================
\section{Future Improvements}
% =========================================================
JWT-based authentication with real role checks; fine payment records;
pagination and Atlas Search; a replica-set deployment enabling true
transactions; generating the OpenAPI spec from code (or tests asserting
spec/route parity); Dockerised deployment; notification preferences;
borrow-history export.

% =========================================================
\section{Conclusion}
% =========================================================
The Smart Library \& Fine Management System delivers five working modules on
the required MERN + Tailwind stack, applying the Incremental Model, three
design patterns in genuine production paths (Singleton, Facade, Strategy),
TDD with recorded Red--Green evidence, mocked unit tests, Supertest
integration tests against an isolated database, a passing Selenium
end-to-end workflow, a GitHub Actions pipeline and hand-authored Swagger
documentation. The engineering artefacts — requirements, diagrams, tests,
CI and documentation — describe the same system that the code implements,
with limitations stated honestly rather than glossed over.

% =========================================================
\section{References}
% =========================================================
\begin{enumerate}[leftmargin=1.6em]\small
  \item Express.js Team, ``Express 4.x Guide,'' \url{https://expressjs.com/en/guide/}, accessed 2026.
  \item MongoDB, Inc., ``Mongoose ODM Documentation,'' \url{https://mongoosejs.com/docs/}, accessed 2026.
  \item MongoDB, Inc., ``MongoDB Manual: Documents and Schema Design,'' \url{https://www.mongodb.com/docs/manual/}, accessed 2026.
  \item React Team, ``React Documentation,'' \url{https://react.dev/}, accessed 2026.
  \item Tailwind Labs, ``Tailwind CSS v3 Documentation,'' \url{https://v3.tailwindcss.com/docs}, accessed 2026.
  \item Jest Team, ``Jest Testing Framework,'' \url{https://jestjs.io/docs/}, accessed 2026.
  \item Supertest contributors, ``supertest,'' \url{https://github.com/ladjs/supertest}, accessed 2026.
  \item Selenium Project, ``Selenium WebDriver Documentation,'' \url{https://www.selenium.dev/documentation/webdriver/}, accessed 2026.
  \item OpenAPI Initiative, ``OpenAPI Specification 3.0.3,'' \url{https://swagger.io/specification/}, accessed 2026.
  \item GitHub, ``GitHub Actions Documentation,'' \url{https://docs.github.com/en/actions}, accessed 2026.
  \item Gamma, E., Helm, R., Johnson, R., Vlissides, J. \emph{Design Patterns: Elements of Reusable Object-Oriented Software}. Addison-Wesley, 1995.
  \item Sommerville, I. \emph{Software Engineering}, 10th ed. Pearson, 2016.
  \item mongodb-memory-server contributors, \url{https://github.com/typegoose/mongodb-memory-server}, accessed 2026.
\end{enumerate}

% =========================================================
\appendix
\section{Key Code Excerpts}
\label{app:code}
% =========================================================
Full source is in the project repository. Additional excerpts: the complete
Singleton (\texttt{server/src/database/DatabaseManager.js}), the heavy-overdue
decorator (\texttt{server/src/strategies/OverdueHeavyFineStrategy.js}) and
the atomic stock decrement (\texttt{server/src/repositories/BookRepository.js}).

\begin{lstlisting}[caption={Heavy-overdue decorator (core)},label=lst:heavy]
calculateFine(overdueDays) {
  const days = Number(overdueDays);
  if (!Number.isInteger(days) || days < 0) {
    throw new Error('overdueDays must be a non-negative whole number');
  }
  if (days <= this.threshold) {
    return this.baseStrategy.calculateFine(days);
  }
  const basePart = this.baseStrategy.calculateFine(this.threshold);
  const extraDays = days - this.threshold;
  return basePart + extraDays * this.heavyRate;
}
\end{lstlisting}

\begin{lstlisting}[caption={Atomic availability decrement},label=lst:atomic]
async decrementAvailable(id) {
  return this.model.findOneAndUpdate(
    { _id: id, availableCopies: { $gt: 0 } },
    { $inc: { availableCopies: -1 } },
    { new: true }
  );
}
\end{lstlisting}

% =========================================================
\section{Test Results}
\label{app:tests}
% =========================================================
\textbf{Unit + integration run (recorded):} 4 suites, 64 tests, all passing
(\texttt{npm run test}).
\placeholder{re-run at submission and paste your terminal output / screenshot.}

\subsection{Coverage Summary}
\label{app:coverage}
Recorded with \texttt{npm run test:coverage} on the development machine:

\begin{table}[H]
\centering
\caption{Jest coverage (actual recorded run)}
\small
\begin{tabular}{lrrrr}
\toprule
File & \% Stmts & \% Branch & \% Funcs & \% Lines \\
\midrule
All files & 89.77 & 73.33 & 91.01 & 90.88 \\
src/database/DatabaseManager.js & 82.60 & 80.00 & 80.00 & 82.60 \\
src/facades/BorrowFacade.js & 95.08 & 82.85 & 100 & 95.08 \\
src/middleware/errorHandler.js & 35.29 & 7.14 & 75.00 & 33.33 \\
src/middleware/validate.js & 92.00 & 90.90 & 100 & 91.66 \\
src/models/* & 96.66 & 50.00 & 100 & 96.66 \\
src/repositories/* & 94.64 & 70.58 & 95.45 & 96.22 \\
src/services/FineService.js & 95.65 & 90.90 & 100 & 95.45 \\
src/services/AuditService.js & 85.71 & 33.33 & 100 & 85.71 \\
src/strategies/* & 93.75 & 90.47 & 83.33 & 93.75 \\
src/utils/* & 100 & 100 & 100 & 100 \\
\bottomrule
\end{tabular}
\end{table}

Interpretation: statement/line coverage is high for business logic; the main
gap is untested branches of the centralised error handler (e.g.\ specific
Mongoose error shapes). Identified improvement: add unit tests that feed
\texttt{ValidationError}, \texttt{CastError} and duplicate-key errors into
\texttt{errorHandler}.

\subsection{TDD Red--Green Record}
\label{app:tdd}
Red run (unimplemented strategy): \texttt{Tests: 2 failed, 16 skipped,
3 passed, 21 total} — failures: ``one overdue day $\rightarrow$ 10 BDT'' and
``five overdue days $\rightarrow$ 50 BDT''. Green run (implemented):
\texttt{Tests: 16 skipped, 5 passed, 21 total}. Full logs:
\texttt{docs/evidence/tdd-red-run.txt}, \texttt{tdd-green-run.txt}.

\subsection{Selenium E2E}
\label{app:e2e}
Recorded output (development machine):

\begin{lstlisting}[language={},caption={e2e run (recorded)}]
STEP 1 OK — dashboard loaded
STEP 2 OK — book added, toast shown: Book "Selenium Test Book ..." added to the catalog.
STEP 3 OK — member registered and visible in table
STEP 4 OK — book issued, toast shown: Book issued to E2E Member. Due date: 2026-11-09
STEP 5 OK — active borrow row visible
STEP 6 OK — book returned, toast shown: Book returned on time. No fine.

E2E TEST PASSED — issue/return workflow works through the UI
\end{lstlisting}

\placeholder{Add a screenshot of your own run and of the headed browser.}

% =========================================================
\section{Git Evidence}
\label{app:git}
% =========================================================
\placeholder{Screenshots: git log graph, your commits, branch list.}

% =========================================================
\section{Pull-Request Evidence}
\label{app:pr}
% =========================================================
\placeholder{Accepted PR with reviews; rejected/revised PR with comments.}

% =========================================================
\section{Merge-Conflict Evidence}
\label{app:conflict}
% =========================================================
\placeholder{Conflict markers screenshot + resolved commit.}

% =========================================================
\section{CI Evidence}
\label{app:ci}
% =========================================================
\placeholder{GitHub Actions successful run screenshot (run number + date).}

% =========================================================
\section{Swagger Evidence}
\label{app:swagger}
% =========================================================
\placeholder{Screenshot of /api-docs with a "Try it out" response.}
Live checks performed during development: \texttt{GET /api-docs $\rightarrow$
200}; \texttt{GET /api-openapi.json} returned OpenAPI 3.0.3 with all
documented paths.

% =========================================================
\section{Learning Log}
\label{app:log}
% =========================================================
\placeholder{Your completed learning log (see docs/learning-log.md).}

% =========================================================
\section{Application Screenshots}
\label{app:screens}
% =========================================================
\placeholder{Dashboard, Members, Books, Issue/Return, Fines, Audit pages
from your running application (add PNGs to figures/ and reference them
here).}

\end{document}
